\documentclass[11pt,a4paper]{article}
\usepackage[T1]{fontenc}
\usepackage[utf8]{inputenc}
\usepackage{lmodern,amsmath,amssymb,amsthm,booktabs,array}
\usepackage[margin=25mm]{geometry}
\usepackage{microtype}
\usepackage[colorlinks=true,linkcolor=black,citecolor=blue,urlcolor=blue]{hyperref}
\hypersetup{pdftitle={Arithmetic-Derivative Pythagorean 7-Tuples: Exact Extension and Homothetic Structure of OEIS A249107},pdfauthor={Carlo Corti}}
\newcommand{\NN}{\mathbb{N}}
\newcommand{\PP}{\mathbb{P}}
\newcommand{\D}{\mathrm{D}}
\newcommand{\Lrate}{\mathrm{L}}
\newtheorem{theorem}{Theorem}[section]
\newtheorem{proposition}[theorem]{Proposition}
\newtheorem{lemma}[theorem]{Lemma}
\newtheorem{corollary}[theorem]{Corollary}
\theoremstyle{definition}
\newtheorem{definition}[theorem]{Definition}
\theoremstyle{remark}
\newtheorem{remark}[theorem]{Remark}
\title{Arithmetic-Derivative Pythagorean 7-Tuples:\\Exact Extension and Homothetic Structure of OEIS A249107}
\author{Carlo Corti\\Independent researcher\\\texttt{carlo.corti@outlook.com}}
\date{October 2026}

\begin{document}
\maketitle
\begin{abstract}
OEIS A249107 lists the integers whose squares, together with those of their
first five arithmetic derivatives, sum to a square. Relative to the accepted
ten-term prefix, an exact segmented computation adds six consecutive terms:
$1058578125$, $1072281223$, $1774631673$, $2539872501$, $9873543431$, and
$30429212500$. The exclusion claim covers all $30048948369$ integers from
$380264132$ through $30429212500$ and depends on the executed C implementation
and its completed segment records. Independent arithmetic reconstructs all
sixteen positive witnesses; independent negative checks cover finite windows,
not the full added interval. We give a complete integer-scale criterion for
each fixed admissible primitive ordered form, with a residue-core proof for
rational logarithmic rates. Applying it to the sixteen witnesses yields eleven
observed forms: two inherited infinite homothetic classes and nine singleton
classes. Three of the six added terms introduce forms absent from the original
prefix. Exact finite trajectory checks and an elementary growth lemma show
that three of the sixteen derivative trajectories terminate at zero and the
other thirteen diverge. One terminating term satisfies the same square identity
at every depth at least five under the explicit zero convention. Finally,
$30429212500<a(17)\le205397184375$; the upper witness is not assigned an index.
\end{abstract}

\noindent\textbf{2020 Mathematics Subject Classification.}
Primary 11A25; Secondary 11D09, 11Y16, 11Y55.

\smallskip
\noindent\textbf{Key words and phrases.}
Arithmetic derivative, Pythagorean tuple, OEIS A249107, exact enumeration,
primitive ordered form, homothetic class, integer scale.

\section{Definition, baseline, and contribution}\label{sec:intro}

The arithmetic derivative on the nonnegative integers is defined by
\begin{equation}\label{eq:derivative}
 \D(0)=\D(1)=0,\qquad
 \D(n)=\sum_{p^e\parallel n}e\frac{n}{p}\quad(n\ge2).
\end{equation}
Here $p$ is prime and $p^e\parallel n$ means that $e$ is its exact exponent.
It satisfies the product rule
\begin{equation}\label{eq:product}
 \D(ab)=a\D(b)+b\D(a)
\end{equation}
without a coprimality hypothesis. Indeed, prime exponents add in a product,
which proves the identity directly from~\eqref{eq:derivative}; zero factors
are covered by $\D(0)=0$. It is not additive: $\D(2)+\D(3)=2$, whereas
$\D(5)=1$. Basic properties, logarithmic equations, and higher iterates are
treated by Ufnarovski and \AA hlander~\cite{ufnarovski} and Kovi\v{c}~\cite{kovic}.
The survey of Tossavainen and coauthors~\cite{overview} provides further
background; our reading used its accepted author manuscript.

Write $\D^0(n)=n$ and $\D^{j+1}(n)=\D(\D^j(n))$; the superscript denotes
composition, not a power of the value. For $h\ge1$, put
\begin{equation}\label{eq:Qh}
 Q_h(n)=\sum_{j=0}^{h}\bigl(\D^j(n)\bigr)^2,
 \qquad S_h=\{n\ge2:Q_h(n)\text{ is an integer square}\}.
\end{equation}
We study the increasing enumeration $a(1),a(2),\ldots$ of $S_5$,
with offset~1. Its witness is the ordered seven-coordinate tuple
\begin{equation}\label{eq:witness}
 W(n)=\bigl(n,\D(n),\D^2(n),\D^3(n),\D^4(n),\D^5(n),R\bigr),
 \qquad R^2=Q_5(n),\quad R>0.
\end{equation}
No primitivity, ordering of magnitudes, common-divisor condition, or
positivity condition on later derivatives is imposed. All derivative
coordinates are nonnegative, and a zero remains zero thereafter.

Paolo P. Lava introduced A249107 on October 21, 2014. The entry credits Ray
Chandler with terms 5--10 in extensions dated December 2016 and January
2017~\cite{sequence}. The accepted baseline used here has ten terms, ending
at $a(10)=380264131$. A direct read of the entry and its b-file on
October 21, 2026 still gives that ten-term baseline. We accept its indexing
and gap completeness as prior data; we have not repeated a census of the
whole domain below it. The Maple call printed in the entry is an algorithm
example, not a retained certificate of an additional negative interval.

This work separates three contributions. First, a completed census beyond
the accepted endpoint assigns six additional consecutive indices. Second,
normalization and a complete fixed-form scale criterion describe exactly
the homothetic classes represented by the sixteen witnesses. Third, finite
derivative chains provide rigorous trajectory information and depth
comparisons for those witnesses. The two specialized seeds used below and
their scaling mechanism were already present in the project analyses before
the census. The general fixed-form argument follows the author's earlier
studies of A230543 and A249105~\cite{corti230543,corti249105}; it is proved here
at depth five, including the integer-positivity restriction on scales.
The exact new ranks, three additional observed forms, and the finite
trajectory analysis must be distinguished from those inherited tools.
No absolute priority claim is made for the specialized seeds or a general
classification of all admissible forms.

\section{Complete scales of one admissible form}\label{sec:scales}

For a positive integer $s$, its logarithmic derivative is
\begin{equation}\label{eq:L}
 \Lrate(s)=\frac{\D(s)}{s}=\sum_{p^e\parallel s}\frac{e}{p},
 \qquad \Lrate(1)=0.
\end{equation}
In particular $\Lrate(st)=\Lrate(s)+\Lrate(t)$, including when $s,t$ have
common prime factors. The rational logarithmic equation and its squarefree
denominator restriction are classical~\cite[Section 8]{ufnarovski}.
Kovi\v{c}'s Proposition~2 classifies the integer-rate case~\cite{kovic}.

\begin{definition}\label{def:form}
For a witness $W$, let $g$ be the gcd of all seven coordinates, including
its root, and let
\[
 F=W/g=(u_0,u_1,u_2,u_3,u_4,u_5,r),\qquad u_0=m.
\]
This is its primitive ordered form. A nonnegative integer tuple of this
shape, with $m,r>0$ and gcd~1, is admissible if some positive integer
multiple is a witness with input at least~2. Its homothetic class consists
of all such witnesses $sF$.
\end{definition}

The order of coordinates is retained. A primitive form need not itself
be a derivative chain, since dividing a chain by $g$ does not commute
with arithmetic differentiation.

\begin{theorem}[Fixed-form criterion]\label{thm:criterion}
Let $F$ be an admissible primitive ordered form and define
\begin{equation}\label{eq:lambda}
 \lambda=\frac{u_1-\D(m)}{m}.
\end{equation}
Then $\lambda\ge0$, and
\begin{equation}\label{eq:shift}
 u_{j+1}=\D(u_j)+\lambda u_j\quad(0\le j<5),
 \qquad \sum_{j=0}^{5}u_j^2=r^2.
\end{equation}
For every positive integer $s$ with $sm\ge2$, the tuple $sF$ is a witness
if and only if $\Lrate(s)=\lambda$. No noninteger positive rational
multiplier makes all seven coordinates of $F$ integral.
\end{theorem}
\begin{proof}
Choose a witnessing scale $t$. From $\D(tm)=tu_1$ and the product rule,
$\D(m)+m\Lrate(t)=u_1$, so $\lambda=\Lrate(t)\ge0$. Each remaining
transition $\D(tu_j)=tu_{j+1}$ gives the corresponding equality
in~\eqref{eq:shift}. If $u_j=0$, this equality forces $u_{j+1}=0$,
so the shifted transition also holds at zero. The geometric identity
follows by dividing the original square identity by $t^2$.

For a proposed $s$, the first derivative equality $\D(sm)=su_1$ is
equivalent to $\Lrate(s)=\lambda$. If it holds, induction gives
\[
 \D(su_j)=s\D(u_j)+u_j\D(s)
          =s\bigl(\D(u_j)+\lambda u_j\bigr)=su_{j+1}
 \quad(0\le j<5).
\]
This remains valid at a zero coordinate. The square identity scales by
$s^2$, proving sufficiency. Finally, if $a/b>0$ is reduced and $(a/b)F$
is integral, then $b$ divides every coordinate of $F$, and hence $b=1$.
\end{proof}

Only the product rule and integer arguments were used. In particular,
there is no binomial expansion of a purported additive operator. For an
unwitnessed proposed form, the square identity, all five transitions,
the scale equation, and the input domain must each be verified.

\begin{theorem}[Residue core and all positive integer scales]\label{thm:residue}
Let $\lambda=A/B\ge0$ be reduced, with $B>0$. If $\Lrate(s)=\lambda$
has a positive integer solution, then $B$ is squarefree. When $B$ is
squarefree, define, for every prime $p\mid B$,
\begin{equation}\label{eq:residues}
 c_p\equiv A(B/p)^{-1}\pmod p,\qquad 1\le c_p<p,
\end{equation}
and put
\begin{equation}\label{eq:core}
 s_0=\prod_{p\mid B}p^{c_p},\qquad
 \alpha_0=\sum_{p\mid B}\frac{c_p}{p},\qquad
 k=\lambda-\alpha_0.
\end{equation}
For $B=1$, use $s_0=1$ and $\alpha_0=0$. Then $k$ is an integer.
There is a positive integer solution if and only if $k\ge0$, in which
case all and only the solutions are
\begin{equation}\label{eq:allscales}
 s=s_0\prod_{j=1}^{k}p_j^{p_j},\qquad p_j\in\PP.
\end{equation}
Prime parameters may repeat and may divide $B$. For an admissible form,
the scales in~\eqref{eq:allscales} are exactly its homothetic class.
\end{theorem}
\begin{proof}
Every $\Lrate(s)=\sum e_p/p$ has reduced denominator dividing the
squarefree product of the primes of $s$. Suppose now that $B$ is
squarefree. The residue representatives exist and are nonzero since
$\gcd(A,B)=1$. Furthermore
\[
 Bk=A-\sum_{p\mid B}c_p(B/p)
\]
is an integer divisible by every prime of $B$, by~\eqref{eq:residues}.
Since $B$ is squarefree, $B$ divides this integer, proving $k\in\mathbb Z$.
For $B=1$ the assertion is immediate.

For any solution $s=\prod p^{e_p}$, let $M$ be the squarefree product of
the primes dividing $sB$. Multiplication of the rate equation by $M$ gives
\begin{equation}\label{eq:modscale}
 \sum_{p\mid s}e_p M/p=A M/B.
\end{equation}
If $p\mid B$ were absent from $s$, reduction modulo $p$ would give zero
on the left and a nonzero number on the right. Thus $p\mid s$. For such
a prime, reduction and cancellation of the unit $M/B$ modulo $p$ give
$e_p(B/p)\equiv A\pmod p$. For a prime of $s$ outside $B$, reduction
instead gives $e_p\equiv0\pmod p$. Consequently
\[
 e_p=c_p+ph_p\ (p\mid B),\qquad
 e_p=ph_p\ (p\nmid B),\qquad h_p\in\mathbb Z_{\ge0}.
\]
Substitution yields $\lambda=\alpha_0+\sum_p h_p$, so $k\ge0$ and
listing each prime $h_p$ times gives~\eqref{eq:allscales}. Conversely,
every displayed product has rate $\alpha_0+k=\lambda$. For $k\ge0$,
choosing any $k$ primes proves existence. The fixed-form conclusion is
Theorem~\ref{thm:criterion}; the input restriction causes no exception for
an admissible form. The only possible failure of $sm\ge2$ would be $s=m=1$.
Since the proposed scale satisfies $\Lrate(s)=\lambda$, it would give
$\lambda=0$ and, by~\eqref{eq:shift}, $u_1=\cdots=u_5=0$.
Any witnessing scale $t$ would then also satisfy $\Lrate(t)=0$, hence
$t=1$ by~\eqref{eq:L}. Its input would be $tm=1$, contradicting
admissibility.
\end{proof}

\begin{corollary}\label{cor:cardinality}
An admissible fixed form has a singleton homothetic class precisely when
$k=0$, and an infinite class when $k\ge1$. At an integer rate $\lambda$,
its scales are products of exactly $\lambda$ factors $p^p$, with repetition.
\end{corollary}
\begin{proof}
For $k=0$ the sole scale is $s_0$. For $k\ge1$, fix $k-1$ prime parameters
and let the remaining prime increase. This gives infinitely many distinct
scales and first coordinates. For an integer rate, $B=1$ and $k=\lambda$.
\end{proof}

The denominator condition alone is insufficient in the integer domain.
For $\lambda=1/6$, one has $c_2=1$, $c_3=2$, $s_0=18$,
$\alpha_0=7/6$, and $k=-1$, so there is no positive integer scale.
In contrast, rate $1/2$ has the sole scale~2. This last example concerns
the rate equation alone, not an exhibited A249107 form. Thus ``singleton''
means $k=0$, not necessarily $\lambda=0$. These results classify the
scales of a fixed form; they do not classify all admissible forms.

\section{The inherited families and the eleven observed forms}\label{sec:forms}

Two primitive ordered seed tuples present before the campaign are
\begin{align}
 F_1={}&(338745,550369,550370,947809,954003,1272007,2036425),
       \label{eq:F1}\\
 F_2={}&(2434337,5075913,11896417,25623221,53701605,\notag\\
       &\hspace{33mm}153944646,165567453).\label{eq:F2}
\end{align}
Their rates are respectively 1 and 2. In particular, $\D(338745)=211624$
and $\D(2434337)=207239$; neither tuple is an unshifted derivative chain.
Table~\ref{tab:transitions} gives the five exact derivative evaluations
needed to check the shifted transitions. Direct summation of the six
squares gives respectively $2036425^2$ and $165567453^2$; the gcd of
each complete seven-tuple is~1.

\begin{table}[htbp]
\centering\small
\begin{tabular}{crrrrr}
\toprule
Form & $\D(u_0)$ & $\D(u_1)$ & $\D(u_2)$ & $\D(u_3)$ & $\D(u_4)$\\
\midrule
$F_1$ &211624&1&397439&6194&318004\\
$F_2$ &207239&1744591&1830387&2455163&46541436\\
\bottomrule
\end{tabular}
\caption{Exact transition checks for the inherited primitive forms.
Adding $\lambda u_j$, with $\lambda=1$ or 2 respectively, gives $u_{j+1}$
at each of the five positions.}\label{tab:transitions}
\end{table}

\begin{proposition}\label{prop:families}
The complete homothetic classes of $F_1$ and $F_2$ have first coordinates
\begin{equation}\label{eq:families}
 n=338745p^p,\qquad n=2434337p^pq^q,
 \qquad p,q\in\PP,
\end{equation}
respectively. Both are infinite subsets of A249107. Repeated parameters
and primes dividing the seed are allowed.
\end{proposition}
\begin{proof}
Equation~\eqref{eq:derivative} gives $\D(p^p)=p^p$. Additivity of
$\Lrate$ on products gives $\Lrate(p^pq^q)=2$, also for $p=q$.
The checked transitions and square identities make the scaled tuples
witnesses by the product-rule induction of Theorem~\ref{thm:criterion}.
Theorem~\ref{thm:residue} proves that these are all their scales.
Infinitude follows from Corollary~\ref{cor:cardinality}.
\end{proof}

\begin{table}[htbp]
\centering\small
\begin{tabular}{clcc}
\toprule
Form & Indices represented & Observed scales $g$ & $(\lambda,k)$\\
\midrule
$F_1$ & $5,6,11$ & $4,27,3125$ & $(1,1)$\\
$F_2$ & $7,9,13,16$ & $16,108,729,12500$ & $(2,2)$\\
$T_1$ & $1$ & $1$ & $(0,0)$\\
$T_2$ & $2$ & $1$ & $(0,0)$\\
$T_3$ & $3$ & $1$ & $(0,0)$\\
$T_4$ & $4$ & $1$ & $(0,0)$\\
$T_8$ & $8$ & $1$ & $(0,0)$\\
$T_{10}$ & $10$ & $1$ & $(0,0)$\\
$T_{12}$ & $12$ & $1$ & $(0,0)$\\
$T_{14}$ & $14$ & $1$ & $(0,0)$\\
$T_{15}$ & $15$ & $1$ & $(0,0)$\\
\bottomrule
\end{tabular}
\caption{All primitive ordered forms observed in the sixteen-term prefix.
The full coordinates of $F_1,F_2$ are~\eqref{eq:F1}--\eqref{eq:F2};
$T_i$ is exactly the seven-coordinate witness in row $i$ of
Table~\ref{tab:witnesses}. Every residue core $s_0$ in this table is~1.
Only the two displayed classes, rather than all possible infinite classes,
are being enumerated.}\label{tab:forms}
\end{table}

\begin{proposition}[Classification of the observed forms]\label{prop:observed}
The sixteen witnesses in Table~\ref{tab:witnesses} have exactly the eleven
primitive ordered forms in Table~\ref{tab:forms}. The nine $T_i$ have
singleton homothetic classes over the entire positive integer scale domain.
\end{proposition}
\begin{proof}
Take the gcd of all seven coordinates in each row. The nonunit gcds are
precisely the seven scales shown for $F_1,F_2$. The other nine rows have
gcd~1 and rate zero, and hence are their own primitive derivative witnesses.
Their first coordinates are distinct, and differ from those of $F_1,F_2$,
so the eleven tuples are distinct. All derivative and square identities
were reconstructed using exhaustive trial division and integer square roots.
For each $T_i$, $\Lrate(s)=0$ forces $s=1$ by~\eqref{eq:L}, proving the
global singleton claim for that fixed form.
\end{proof}

The original ten rows represent eight forms, including six singletons.
The added indices 12, 14, and 15 introduce three forms absent from that
sample; indices 11, 13, and 16 enlarge the represented scales of the
inherited families. Specifically, they correspond to $F_1(5)$,
$F_2(3,3)$, and $F_2(2,5)$. Their positive values were available from
the family constructions before the missing intervals were enumerated.
The census supplies their exact indices. All nine singleton first
coordinates in this sample are squarefree; the other seven first
coordinates are not. This is a finite observation, not a general equivalence
between squarefreeness and singleton homothetic classes.

\section{Derivative trajectories and the effect of changing depth}\label{sec:dynamics}

The square identity concerns a finite initial trajectory. Its later behavior
can nevertheless be decided for every term in the present prefix, using
only three additional derivatives and the following historical growth
mechanism~\cite[Theorem 4]{ufnarovski}.

\begin{lemma}\label{lem:growth}
If an iterate of $n$ is $p^pb$ for a prime $p$ and an integer $b>1$,
then $\D^j(n)\longrightarrow\infty$ as $j\longrightarrow\infty$.
\end{lemma}
\begin{proof}
The product rule gives $\D(p^pb)=p^p(b+\D(b))$.
Starting at this iterate, put $b_{t+1}=b_t+\D(b_t)$.
Every $b_t>1$ has a nonempty prime factorization, so
\eqref{eq:derivative} gives $\D(b_t)\ge1$ and $b_t\ge b_0+t$.
Thus the trajectory is unbounded and tends to infinity. No conjecture
about arbitrary arithmetic-derivative trajectories is needed.
\end{proof}

\begin{table}[htbp]
\centering\small
\begin{tabular}{crrrl}
\toprule
$i$ & $\D^6(a(i))$ & $\D^7(a(i))$ & $\D^8(a(i))$ & Eventual behavior\\
\midrule
1&12032&48384&248832&diverges; $4\mid\D(a(1))=168$\\
2&10260&23112&57996&diverges; $4\mid\D(a(2))=216$\\
3&4579200&22524480&104443776&diverges; $4\mid\D(a(3))=24544$\\
4&1202048&4207296&14275072&diverges; $4\mid\D(a(4))=7484$\\
8&1&0&0&first zero at depth 7\\
10&67928168&117772180&143351236&diverges; $4\mid\D(a(10))=29251100$\\
12&0&0&0&first zero at depth 6\\
14&201791&1&0&first zero at depth 8\\
15&4208230&3269601&2285484&diverges; $4\mid\D^8(a(15))$\\
\bottomrule
\end{tabular}
\caption{Exact continuation of the nine singleton witnesses. Every displayed
multiple of 4 used for divergence is greater than 4. The seven represented
members of $F_1,F_2$ already contain a proper factor $p^p$ at depth zero.}
\label{tab:dynamics}
\end{table}

\begin{proposition}[Complete trajectory decision for the finite prefix]\label{prop:dynamics}
For $i=8,12,14$, the trajectory of $a(i)$ first reaches zero at depths
$7,6,8$, respectively. For every other $1\le i\le16$, it tends to infinity.
\end{proposition}
\begin{proof}
The exact continuations in Table~\ref{tab:dynamics}, together with the
positive coordinates of Table~\ref{tab:witnesses}, prove the three first-zero
claims. The six other singleton rows have a displayed iterate equal to
$4b$ with $b>1$, so Lemma~\ref{lem:growth} applies with $p=2$.
For $F_1$, every scale $p^p$ is a proper divisor of its first coordinate,
since $338745>1$. For $F_2$, one factor $p^p$ leaves the multiplier
$2434337q^q>1$. The lemma therefore applies to every member of both
families, including repeated or shared primes.
\end{proof}

This proposition decides sixteen particular trajectories, and in addition
the trajectories of all members of the two proved families. It does not
settle the general dynamics of $\D$, or prove that no other behavior can
occur outside these sets. A singleton homothetic class can thus contain
either a terminating or a divergent trajectory.

\begin{proposition}[A persistent square identity]\label{prop:tail}
Let $N=1072281223=a(12)$ and $R=1076009463$. Then
\begin{equation}\label{eq:persistence}
 \D^5(N)=1,\quad \D^j(N)=0\ (j\ge6),\quad Q_h(N)=R^2\ (h\ge5).
\end{equation}
For the zero convention in~\eqref{eq:derivative}, the sixteen-term
sample has no member of $S_h$ for $1\le h\le4$, and its only member
of $S_h$ for each $h=6,7,8$ is $N$.
\end{proposition}
\begin{proof}
The witness in row 12 has fifth derivative~1. The definitions give its
zero tail, and adding zero squares proves~\eqref{eq:persistence} at every
later depth. For the finite assertions, reconstruct the sixteen chains
through depth 8 and test each partial sum using $t=\lfloor\sqrt{Q_h}\rfloor$
and $t^2=Q_h$. No input interval outside these sixteen values is searched.
\end{proof}

The other two terminating cases illustrate why termination alone is not
enough. For $a(8)$, the stabilized sum from depth 6 onward is
$53881926^2+1$, strictly between consecutive squares. For $a(14)$ it
stabilizes from depth 7 onward at
$2705738777^2+201791^2+1$, also nonsquare by exact integer-root testing.
Different square-sum depths are not generally nested: the fixed point
$n=4$ satisfies $Q_3(4)=64$ and $Q_5(4)=96$, while $4031\in S_5$
and $4031\notin S_3$ by its exact partial sum.

\section{Exact enumeration and numerical domain}\label{sec:algorithm}

\subsection{Arithmetic kernel and coverage}
The production program is a specialized C17 implementation for $Q_5$,
with GCC's unsigned 128-bit extension, OpenMP, GMP, and OpenSSL libcrypto.
It enumerates every integer in a configured inclusive interval, using
segmentation for arithmetic and recovery rather than a restriction to
known families.

For a block $[L,H]$, a prime sieve supplies all primes through
$\lfloor\sqrt H\rfloor$. At each multiple $n$ of $p$, divide its residual
by $p$ until no longer divisible, and add $e(n/p)$ to the first derivative,
where $n$ is the original input. Any residual greater than~1 is prime:
a composite residual would have a prime divisor at most $\sqrt H$ that
the sieve had already removed. Adding $n/\text{residual}$ completes
$\D(n)$ exactly.

For nonprime inputs, four further factorizations compute
$\D^2(n),\ldots,\D^5(n)$. The word factorizer first tests primality using
strong probable-prime tests with bases
\begin{equation}\label{eq:bases}
 2,325,9375,28178,450775,9780504,1795265022.
\end{equation}
The deterministic coverage for inputs below $2^{64}$ is the finite-domain
record attributed to Jim Sinclair in the Miller--Rabin catalogue
\cite{bases}; the machine-word setting is also discussed by Fori\v{s}ek
and Jan\v{c}ina~\cite{forisek}. This external finite-domain result is a
premise of the computation, not a new exhaustive primality verification
in this project. Bases are reduced modulo the tested integer; a zero
residue is skipped while the remaining bases are still tested.

Composite operands are split by fixed-seed Pollard-rho iterations,
using the factor-finding principle of~\cite{pollard}. The source uses
$x\mapsto x^2+c$ with $1\le c\le32$, at most $100000$ steps per seed,
and gcd comparisons of one-step and two-step trajectories. Failure of
the bounded fast path leads to exhaustive odd trial division, after
small primes through~97 have been tried. The trial condition $p\le z/p$
avoids a squaring overflow. Every recursive split is a proper divisor,
so the operands strictly decrease. At prime leaves the program checks
the reconstructed product and accumulates exact multiplicities in
\eqref{eq:derivative}. The finite fallback and proper descent establish
termination in the absence of interruption or resource failure, although
they supply no useful worst-case runtime guarantee for every operand.

The only early mathematical rejection is $\D(n)=1$: this characterizes
prime original inputs, for which $Q_5(n)=n^2+1$ lies strictly between
$n^2$ and $(n+1)^2$. Subsequent ones and zeros are evaluated normally.
Indeed a composite $n=ab$, with $a,b\ge2$, has
$\D(n)=a\D(b)+b\D(a)\ge a+b\ge4$.
There is no family filter, floating-point square decision, or probable-prime
acceptance beyond the stated finite deterministic primality domain.
Each of the six coordinates is imported into GMP before squaring and
summation. GMP's exact square predicate decides membership; an integer
square root is then obtained for a positive result~\cite{gmp}.

\subsection{Integer bounds}
Let $\Omega(z)$ count prime factors with multiplicity. For $z\ge2$,
\begin{equation}\label{eq:width}
 \D(z)\le z\Omega(z)/2\le z\lfloor\log_2 z\rfloor/2.
\end{equation}
The first inequality follows because each reciprocal prime is at most
$1/2$; the second follows from $2^{\Omega(z)}\le z$.
This is the usual arithmetic-derivative bound~\cite{ufnarovski}, in a
convenient integer form. The function
$U\mapsto\lfloor U\lfloor\log_2 U\rfloor/2\rfloor$ is nondecreasing
for integer $U\ge2$: both nonnegative factors $U$ and
$\lfloor\log_2 U\rfloor$ are nondecreasing, as is taking the floor. Iterating it at the supported input cap gives
Table~\ref{tab:bounds}, with 0 and 1 handled separately.

\begin{table}[htbp]
\centering\small
\begin{tabular}{crl}
\toprule
Quantity & Universal upper bound at the source cap & Representation\\
\midrule
$n$ &40000000000000&unsigned word\\
$\D(n)$ &900000000000000&unsigned word\\
$\D^2(n)$ &22050000000000000&unsigned word\\
$\D^3(n)$ &595350000000000000&unsigned word\\
$\D^4(n)$ &17562825000000000000&unsigned word; last factored operand\\
$\D^5(n)$ &553228987500000000000&unsigned 128-bit\\
$\lfloor\sqrt{Q_5(n)}\rfloor$ &553508012914801786476&unsigned 128-bit\\
\bottomrule
\end{tabular}
\caption{Conservative representation bounds, not observed coordinate maxima.
The supported cap $4\cdot10^{13}$ is distinct from the completed census
endpoint $30429212500$.}\label{tab:bounds}
\end{table}

The sum of the six squared bounds is
\[
 Q_5(n)\le306371120360892381250000000000000000000000,
\]
which requires 138 bits. Thus even an unsigned 128-bit sum is not
sufficient at the supported cap; GMP is necessary for this implementation's
full domain. The last factored operand fits below $2^{64}$, so no
larger-integer factorization theorem is assumed. Word modular products
use 128-bit intermediates. Prime multiplicity is at most~63 for these
word operands, within the 64-slot factor array. Decimal conversion,
allocation growth, loop endpoints, and narrowing conversions are checked.
These bounds establish representation capacity; they do not assert the
practical feasibility of a complete census to the cap.

\subsection{Parallel work and completed-block recovery}
The first-derivative sieve is serial within a block. The remaining
per-input work uses OpenMP dynamic scheduling in chunks of~64; each input
owns one result flag. After the join, an incomplete flag prevents the
block from being committed. Positive rows are replayed and appended in
input order serially. The frontier advances only after every input in
the block has completed; an interrupted block and its tentative hits
are discarded and recomputed on resumption. There is no resumption
inside a single factorization or input.

For $N_I$ inputs and block size $\beta$, let $K=\lceil N_I/\beta\rceil$.
If $G$ is the maximum block count per operational segment, set
$\sigma=\lceil K/G\rceil$. Segment $j$, for $0\le j<\sigma$, owns block indices
\[
 \lfloor jK/\sigma\rfloor,\ldots,\lfloor(j+1)K/\sigma\rfloor-1.
\]
These intervals are disjoint and cover all $K$ blocks; adjacent endpoints
coincide, and their block counts differ by at most one. Only the final
input block is clipped to the configured endpoint. The campaign used
$\beta=262144$, eight workers, and four immutable configurations.

Persistent states use temporary-file replacement with file and directory
synchronization, internal digests, and a separate seal for a complete
segment. Immutable configuration and source/build/executable identity
checks bind the run records to the executed program. These checks detect
accidental corruption and inconsistent coverage, rather than protect
against a trusted owner rewriting all evidence. Sealed segment states
are authoritative; aggregate tables are checked against them. A successful
single segment is not completion of the whole configured target.
Mathematical coverage requires every configured segment to be complete,
not merely an exit status or the largest input assigned to a worker.

\section{Computational result and validation}\label{sec:results}

\begin{table}[htbp]
\centering\scriptsize
\setlength{\tabcolsep}{3pt}
\begin{tabular}{rrrrrrrr}
\toprule
$i$&$a(i)$&$\D(a(i))$&$\D^2(a(i))$&$\D^3(a(i))$&$\D^4(a(i))$&$\D^5(a(i))$&$R$\\
\midrule
1&4031&168&332&336&832&2560&4873\\
2&10823&216&540&1188&2484&5076&12281\\
3&416959&24544&63664&130464&457488&1526256&1653567\\
4&496939&7484&7488&28032&107840&345408&615453\\
5&1354980&2201476&2201480&3791236&3816012&5088028&8145700\\
6&9146115&14859963&14859990&25590843&25758081&34344189&54983475\\
7&38949392&81214608&190342672&409971536&859225680&2463114336&2649079248\\
8&44472866&22258485&19398957&6466322&3262101&1100831&53881926\\
9&262908396&548198604&1284813036&2767307868&5799773340&16626021768&17881284924\\
10&380264131&29251100&41359540&49631468&49948220&67928164&395903259\\
\midrule
11&1058578125&1719903125&1719906250&2961903125&2981259375&3975021875&6363828125\\
12&1072281223&85244735&27198937&1431542&920291&1&1076009463\\
13&1774631673&3700340577&8672487993&18679328109&39148470045&112225646934&120698673237\\
14&2539872501&928555631&79073578&39536791&5891811&1986491&2705738777\\
15&9873543431&747786611&127012442&64347665&15088563&12624681&9902863451\\
16&30429212500&63448912500&148705212500&320290262500&671270062500&1924308075000&2069593162500\\
\bottomrule
\end{tabular}
\caption{The complete locally consolidated sixteen-term prefix and all
ordered witnesses. Each row has five derivative transitions and satisfies
$a(i)^2+\D(a(i))^2+\cdots+\D^5(a(i))^2=R^2$. The horizontal break separates
the inherited ten-term baseline from the six additional indexed terms.
No coordinate is sorted or suppressed.}\label{tab:witnesses}
\end{table}

\begin{table}[htbp]
\centering\small
\begin{tabular}{crrrr}
\toprule
Stage & Inclusive input interval & Segments & Positives & Scan seconds\\
\midrule
1&$380264132\ldots1058578125$&1&1&584.88\\
2&$1058578126\ldots1774631673$&1&2&666.78\\
3&$1774631674\ldots10000000000$&13&2&9243.99\\
4&$10000000001\ldots30429212500$&35&1&26457.62\\
\midrule
Total&$380264132\ldots30429212500$&50&6&36953.27\\
\bottomrule
\end{tabular}
\caption{Completed consecutive campaign intervals. Times are per-stage totals
of the 50 recorded GNU time wall measurements of scanning invocations,
one per segment, shown at their recorded precision. They exclude pauses,
subsequent checks, and other work. Environment and timing scope are stated in
Section~\ref{sec:validation}.}\label{tab:campaign}
\end{table}

\begin{theorem}[Computational extension, with explicit dependencies]\label{thm:census}
Accept the indexed ten-term prefix in~\cite{sequence}. Subject to the
correct execution of the reviewed exact kernel of Section~\ref{sec:algorithm},
including its stated finite-domain primality premise, and the completed
segment evidence in Table~\ref{tab:campaign}, the only members of $S_5$
in the inclusive interval $[380264132,30429212500]$ are
\[
 \begin{gathered}
 1058578125,\quad1072281223,\quad1774631673,\\
 2539872501,\quad9873543431,\quad30429212500.
 \end{gathered}
\]
They are respectively $a(11),\ldots,a(16)$. Their positive witnesses are
independently reconstructed in Table~\ref{tab:witnesses}; completeness
of the new negative census depends on the production execution and seals,
and has not been independently repeated over the full interval.
\end{theorem}
\begin{proof}[Computational proof]
The sieve and exact factor/square kernel test the defining predicate at
every completed input, as described in Section~\ref{sec:algorithm}.
All 50 segment states are complete, with the frontier equal to their
inclusive high endpoint, and agree with their seals and aggregate hit
records. Their within-stage partitions and cross-stage joins have neither
gaps nor overlaps. The union is exactly the interval in the statement.
Its cardinality is
\[
 30429212500-380264132+1=30048948369.
\]
The ordered union of the retained hits contains exactly the six displayed
values, leaving $30048948369-6=30048948363$ negatives in the production
census. The independently checked witnesses prove their membership.
Given the accepted baseline and the production negative evidence, the
six increasing hits receive exactly indices 11--16. File-integrity and
positive checks by themselves would not prove the exclusion claim.
\end{proof}

\begin{corollary}\label{cor:next}
Under the computational dependencies of Theorem~\ref{thm:census},
\begin{equation}\label{eq:next}
 30429212500<a(17)\le205397184375.
\end{equation}
\end{corollary}
\begin{proof}
The strict lower bound follows from complete coverage through $a(16)$.
The family $F_2$ at $(p,q)=(3,5)$ gives
$2434337\cdot3^3\cdot5^5=205397184375$, with witness
\[
 \begin{gathered}
 (205397184375,428280159375,1003760184375,\\
 2161959271875,4531072921875,12989079506250,\\
 13969753846875).
 \end{gathered}
\]
Its five derivative identities and square identity were also checked
directly. It is greater than the completed endpoint, so a next term exists
no later than this positive. No coverage of the intervening gap, or index
for this upper witness, is asserted.
\end{proof}

\subsection{What the validation establishes}\label{sec:validation}
The production run used an AMD Ryzen 9 7940HS with eight physical cores
and eight campaign workers, Linux Mint 22.3 with kernel 6.17.0-42,
GCC 13.3.0, GMP 6.3.0, and OpenSSL 3.0.13. The frozen build uses
\texttt{-O3 -std=c17 -fopenmp}, strict warnings promoted to errors, and
links GMP and libcrypto; no processor-specific tuning flags were added.
The recorded scan time sums 50 GNU time measurements, one per completed
segment. Table~\ref{tab:campaign} aggregates them by stage; the total is
approximately $36953$ seconds, about 10 hours 16 minutes. It is active scan wall time,
not elapsed calendar time; native monotonic nanoseconds have a slightly
different measurement scope. Maximum recorded campaign resident memory
was $10656$ KiB. This measured RSS is distinct from conservative
allocation bounds and from earlier diagnostic RSS measurements.

The archived independent Python checker uses exhaustive trial division and
\texttt{math.isqrt}, without importing the C sieve, Miller--Rabin,
rho, or GMP arithmetic. Each stage invocation checks the ten baseline
positives and that stage's retained hits. A separate finite arithmetic
script, \texttt{finite\_audit.py}, reconstructs all sixteen full witnesses:
80 derivative identities and 16 square identities. It also checks all
normalized transitions, both seeds, the unindexed upper witness, and
48 extra transitions through depth~8, three for each of the sixteen inputs.
Table~\ref{tab:dynamics} displays the 27 continuations for the nine
singleton witnesses. These are finite positive arithmetic and structural
checks, not a second global census.

Independent trial-division enumeration was additionally recorded for one
4096-input window at each end of each of the four stages, giving eight
disjoint windows and $32768$ distinct inputs. They contain three positives and $32765$
negatives. The exact intervals are
\begin{equation}\label{eq:windows}
\begin{array}{rr}
\lbrack380264132,380268227\rbrack, &\lbrack1058574030,1058578125\rbrack,\\
\lbrack1058578126,1058582221\rbrack, &\lbrack1774627578,1774631673\rbrack,\\
\lbrack1774631674,1774635769\rbrack, &\lbrack9999995905,10000000000\rbrack,\\
\lbrack10000000001,10000004096\rbrack, &\lbrack30429208405,30429212500\rbrack.
\end{array}
\end{equation}
Separate development windows contain $5098$ distinct inputs, with four
positives and $5094$ negatives. That count is not added to the post-run
window count: the development and campaign sets have overlaps.

The retained development matrix has sixteen obligation groups. Strict
GCC and Clang 18.1.3 builds, independent small windows and positive
fixtures, one/four/eight-worker agreement, balanced partitions, and
signal/recovery checks exercise different contracts. Numerical tests
include hard factor operands, forced exact fallback, divisors of the
fixed primality bases, cap/refusal boundaries, and genuine wide paths:
one input produces a 68-bit fifth derivative and a 136-bit sum. Synthetic
wide root/elapsed records test serialization and are not presented as
mathematical witnesses. Corrupt or forged records, identity mismatches,
locking, malformed configuration, and injected synchronization/rename/write
failures test fail-closed persistence behavior.

ASan/UBSan checks passed in their recorded bounded scopes with
LeakSanitizer disabled because of its launch-environment limitation.
Qualified OpenMP ThreadSanitizer checks used Clang/libomp/Archer, a silent
synchronized control and a deliberately racy control that was detected,
before testing the candidate at one, four, and eight workers. Static
analysis records include cppcheck and clang-tidy. These tests supplement
the numerical and ownership arguments; neither sanitizer silence nor
thread agreement is a universal proof of absence of defects. External
code reports were adjudicated against the source and these records;
their verdicts do not replace arithmetic evidence or constitute review
of this manuscript.

\subsection{Reproducing a record audit without changing the baseline}\label{sec:reproduce}
The compact publication package retains both canonical tables, the original
checker, and the complete sealed records of all four stages. From its root,
\begin{verbatim}
python3 -B check/verify.py
python3 -B check/verify.py windows
python3 -B check/verify.py native
\end{verbatim}
The default command verifies the exact payload inventory, archive members,
record seals and stage joins; it regenerates both canonical tables and
compares the finite arithmetic output. The \texttt{windows} mode additionally
enumerates the recorded independent windows. The \texttt{native} mode builds
and tests the public source in temporary storage and uses the exact original
executable to audit all four historical stages. These finite checks do not
repeat the global negative census.

The historical checker deliberately binds its ten-term baseline by identity;
the consolidated b-file has sixteen rows. The driver restores the original
reference layout in temporary storage, without changing that baseline or any
historical checker constant. Its historical \texttt{UPPER} constant denotes
$1058578125=a(11)$, not the present bound $205397184375$.
The lossless archives use compact member names and contain explicit mappings,
byte sizes and hashes for restoring original paths. The original source,
build file and executed binary retain their historical identities. The public
build file changes only the two \texttt{check} command paths to use this driver;
its compiler flags and the C source are unchanged. Its new identity and any
new executable are distinct from the historical run identities. Explicit
Python failure checks remain active with optimization enabled. The README
gives dependencies and optional restoration instructions. No PDF is required
to perform these checks.

The two canonical tables have distinct roles. The b-file lists consecutive
index--value pairs through index~16. The documentary a-file contains the
eight columns of Table~\ref{tab:witnesses} and makes every positive witness
directly checkable. It is not a sparse list of later indexed terms beyond
an unresolved gap. The intended release citation is~\cite{project}; its locators are prospective,
as clarified in Data availability below.

\section{Related sequences, literature, and open questions}\label{sec:context}

OEIS A003415 is the arithmetic derivative itself~\cite{derivativeentry}.
The lower-depth square-sum entries A230543, A249105, and A249106 use
depths $2,3,4$ respectively, while A249110 uses depth~8
\cite{depthentries}. A210503 concerns primitive Pythagorean triples at
depth~1~\cite{triples}; its primitivity condition is additional to a
square-sum test. Thus the family constants, positivity policies, widths,
term indices, and validation domains from these neighboring problems
cannot be transferred to A249107. Only the underlying product-rule
reasoning applies after the relevant transitions have been checked.

The persistence in Proposition~\ref{prop:tail} is an unconditional
statement for the explicit $S_h$ convention of this paper. In particular,
it supplies a zero-extended depth-eight witness with the same root.
The current A249110 entry's displayed Maple factorization expression
does not explicitly implement a separate $\D(0)=0$ branch. We have not
executed that Maple program on a zero tail or assigned $N$ an index in
that entry. Identification with an implemented external predicate must
therefore retain the zero-convention qualification. The original A249107
comment about no eight- or nine-tuple examples up to $10^8$ has that
different finite domain and does not contradict $N>10^9$.

The entries A258644 and A258645 give fourth and fifth arithmetic
derivatives as functions of their input index; A258651 is the array
$A(n,k)=\D^k(n)$, flattened by antidiagonals~\cite{arrayentries}.
Their values encode coordinates of a chain. A square identity and a
rank in A249107 do not follow from being a value in one of these arrays.
The array's explicit zero and unit conventions agree with those used here.

The theoretical distinction relevant to the scale classification is the
domain of the logarithmic equation. Ufnarovski and \AA hlander study
rational solutions, while Kovi\v{c} treats positive integer solutions at
integer rates~\cite{ufnarovski,kovic}. The complete integer-scale statement
here adds the residue-core nonnegativity condition to the squarefree
denominator requirement, following the author's fixed-form work
\cite{corti230543,corti249105}. Deza and coauthors~\cite{deza} examine
arithmetic constants and rational differential equations, including
multiplicative relations between solutions. These rational constructions
are relevant background, but do not replace the integrality, all-five-
transition, and square-identity conditions for a depth-five witness.
The cited works are used for these specific contributions; no longer
ranked A249107 dataset was admitted from them.

Several questions remain open within this study. What is $a(17)$ inside
the interval~\eqref{eq:next}? Are there infinitely many distinct admissible
primitive ordered forms at depth five, or infinitely many singleton
forms? Does an admissible primitive form of this sequence have a noninteger
rate? The observed eleven forms have rates only $0,1,2$, but this finite
sample and the fixed-form theorem supply no answer for unobserved forms.
Likewise, whether there are infinitely many terminating members of $S_5$
is separate from the three terminating examples proved here. No density
law or global conjecture is proposed on the strength of sixteen terms.
The existence and infinitude of the two displayed classes, in contrast,
are theorems and require no such extrapolation.

\section{Conclusion}\label{sec:conclusion}

Complete added interval coverage assigns six consecutive terms beyond the
accepted ten-term prefix, ending at $a(16)=30429212500$, with the stated
production dependence of the negative census. The positive witnesses,
their primitive forms, and their scale transitions admit short exact
checks independent of that census. Three additional forms occur among
the new terms, while three of the added indices belong to families whose
positive values were already available before enumeration. The resulting
eleven observed forms separate two infinite homothetic classes from nine
singleton classes without exhausting the possible forms of the sequence.
The finite trajectory decision also separates homothetic structure from
dynamics: both divergence and eventual zero occur among singleton forms.
The persistent identity at $a(12)$ illustrates why the convention for
later zero coordinates matters in comparisons across depths. The next
rank remains bounded by a proved family witness, and its exact value
requires additional ordering evidence.

\section*{Data availability}
At the research cutoff, the code, immutable configurations, sealed segment
records, result tables, canonical b-file, documentary witness table,
independent checker and frozen reference inputs, finite manuscript-arithmetic
script and its outputs, historical environment captures, and validation
summaries are retained locally. A public release of this project has not been verified,
and this draft does not claim that these materials are already publicly
accessible. The project repository and archive listed in the References
identify the intended release locations. The GitHub URL identifies the
intended public release repository; the literal Zenodo DOI placeholder
will be replaced after the first publication on Zenodo when the author
supplies the verified designated official locator. The mathematical tables
and proofs needed to read the results are contained in this monolithic source.

\section*{AI use disclosure}
OpenAI Codex assisted with the project analysis, implementation workflow,
local artifact checks, and preparation of this manuscript. The supplied code
review material explicitly reports contributions from Anthropic Claude and
the GLM family; other exact backend identities and versions are not uniformly
documented. The supplied manuscript-review reports identify OpenAI ChatGPT,
Anthropic Claude, GLM, MiniMax, and Qwen; several other reports give no
model identity. These are reported identities, not independently
authenticated backend versions. These reported tools are not coauthors, in accordance with the
responsibility and transparency principles of~\cite{cope}.

\paragraph{Mathematical analysis.}
AI assistance supported the reconstruction of the sequence definition from
OEIS, the product-rule and scale proofs, normalization of the full ordered
witnesses, and investigation of terminating and divergent trajectories.
The mathematical statements are supported by the displayed proofs, primary
source readings, exact witness tables, and separately executed finite
integer calculations.

\paragraph{Algorithm and code.}
Codex assisted with the specialized C implementation and its numerical,
parallel, and recovery contracts. The supplied Claude and GLM contributions
include external code review; they do not establish independent enumeration.
The final computational claims are tied to the executed program, immutable
run records, seals, known-term replay, independent trial-division checker,
and scoped sanitizer and fault-test evidence. Unverified execution statements
in review prose were not substituted for local validation records.

\paragraph{Manuscript.}
Codex assisted in drafting English prose and LaTeX, checking literature
metadata and contributor credit, and internally reviewing claim scope,
chronology, and availability. This revision incorporates pertinent findings from the supplied external
AI manuscript reviews after adjudication against the source and retained
evidence. The reports' execution statements have not been added to the
local validation results. Final numerical claims rely on local artifacts
and validation, not conversational assertions. Carlo Corti retains full
scientific responsibility for the mathematical statements, algorithm,
implementation, numerical results, validation boundaries, and conclusions.

\clearpage
\begin{thebibliography}{99}

\bibitem{sequence}
P. P. Lava, OEIS A249107, Numbers that form a Pythagorean 7-tuple with
their first five arithmetic derivatives, October 21, 2014; extensions by
R. Chandler, December 22--23, 2016 and January 2 and 8, 2017.
\url{https://oeis.org/A249107}. Entry and ten-term b-file consulted
October 21, 2026.

\bibitem{ufnarovski}
V. Ufnarovski and B. \AA hlander, How to Differentiate a Number,
\emph{Journal of Integer Sequences} \textbf{6} (2003), Article 03.3.4.
\url{https://cs.uwaterloo.ca/journals/JIS/VOL6/Ufnarovski/ufnarovski.html}.

\bibitem{kovic}
J. Kovi\v{c}, The Arithmetic Derivative and Antiderivative,
\emph{Journal of Integer Sequences} \textbf{15} (2012), Article 12.3.8.
\url{https://cs.uwaterloo.ca/journals/JIS/VOL15/Kovic/kovic4.html}.

\bibitem{overview}
T. Tossavainen, P. Haukkanen, J. K. Merikoski, and M. Mattila,
We Can Differentiate Numbers, Too, \emph{The College Mathematics Journal}
\textbf{55} (2024), no. 2, 100--108; first published online in 2023.
\url{https://doi.org/10.1080/07468342.2023.2268494}.
Accepted author manuscript consulted; no version-specific equation numbering
from the published article is used.

\bibitem{corti230543}
C. Corti, \emph{An Exact Local Extension and Homothetic Classes for OEIS
A230543}, 2026. Manuscript in the archive series \emph{OEIS A230543:
exact extension and homothetic classes}, concept DOI
\url{https://doi.org/10.5281/zenodo.23311392}.

\bibitem{corti249105}
C. Corti, \emph{A Consecutive Extension and Fixed-Form Families for OEIS
A249105}, 2026. Manuscript in the archive series \emph{A249105: Exact
Extension and Fixed-Form Families}, concept DOI
\url{https://doi.org/10.5281/zenodo.23388128}.

\bibitem{deza}
E. I. Deza, T. N. Kazarikhina, D. Y. Ganusenko, and S. O. Tikhonov,
Some questions related to the arithmetic derivative and its applications,
\emph{Chebyshevskii Sbornik} \textbf{26} (2025), no. 3, 96--112 (in Russian).
\url{https://doi.org/10.22405/2226-8383-2025-26-3-96-112}.

\bibitem{bases}
W. Izykowski, \emph{Deterministic variants of the Miller--Rabin primality
test: SPRP bases records}, Jim Sinclair's seven-base record dated
April 20, 2011. \url{https://miller-rabin.appspot.com/}.
Consulted October 21, 2026.

\bibitem{forisek}
M. Fori\v{s}ek and J. Jan\v{c}ina, Fast Primality Testing for Integers That
Fit into a Machine Word, in \emph{Proceedings of Student Research Forum
Papers and Posters at SOFSEM 2015}, CEUR Workshop Proceedings,
Vol. 1326, 2015, pp. 20--30.
\url{https://ceur-ws.org/Vol-1326/020-Forisek.pdf}.

\bibitem{pollard}
J. M. Pollard, A Monte Carlo method for factorization,
\emph{BIT} \textbf{15} (1975), no. 3, 331--334.
\url{https://doi.org/10.1007/BF01933667}.

\bibitem{gmp}
The GMP Development Team, \emph{GNU MP 6.3.0 Manual}, sections on integer
roots and integer import/export.
\url{https://gmplib.org/manual/Integer-Roots};
\url{https://gmplib.org/manual/Integer-Import-and-Export}.
Consulted October 21, 2026.

\bibitem{derivativeentry}
N. J. A. Sloane and R. K. Guy, OEIS A003415, Arithmetic derivative of $n$.
\url{https://oeis.org/A003415}. Consulted October 21, 2026.

\bibitem{depthentries}
P. P. Lava, OEIS A230543 (October 25, 2013), A249105, A249106, and A249110
(October 21, 2014), arithmetic-derivative square-sum entries at depths
2, 3, 4, and 8. Extensions to A230543 include G. Resta (2013),
R. Chandler (2016--2017), and C. Corti (2026); extensions to A249105 and
A249106 include R. Chandler (2016--2017) and C. Corti (2026).
\url{https://oeis.org/A230543}; \url{https://oeis.org/A249105};
\url{https://oeis.org/A249106}; \url{https://oeis.org/A249110}.
Consulted October 21, 2026.

\bibitem{triples}
P. P. Lava, OEIS A210503, Numbers forming a primitive Pythagorean triple
with their arithmetic derivative, January 25, 2013; table and further
comments by R. Chandler.
\url{https://oeis.org/A210503}. Consulted October 21, 2026.

\bibitem{arrayentries}
A. P. Heinz, OEIS A258644, A258645, and A258651, fourth and fifth
arithmetic derivatives and the array of iterated arithmetic derivatives,
June 6, 2015; Mathematica program for A258651 by J.-F. Alcover, 2017.
\url{https://oeis.org/A258644}; \url{https://oeis.org/A258645};
\url{https://oeis.org/A258651}. Consulted October 21, 2026.

\bibitem{cope}
COPE Council, \emph{COPE position---Authorship and AI---English},
version 1, 2024. \url{https://doi.org/10.24318/cCVRZBms}.
Consulted October 21, 2026.

\bibitem{project}
C. Corti, \emph{Arithmetic-Derivative Pythagorean 7-Tuples: Exact Extension
and Homothetic Structure of OEIS A249107}, GitHub repository and Zenodo archive,
2026. \url{https://github.com/carcorti/A249107};
\url{https://doi.org/10.5281/zenodo.xxxxxxxx}.
Author-designated draft locators for the intended release; public
availability has not been verified.

\end{thebibliography}
\end{document}
